\documentclass[11pt,a4paper]{article}

\usepackage[utf8]{inputenc}
\usepackage[T1]{fontenc}
\usepackage[provide=*,spanish]{babel}
\usepackage{lmodern}
\usepackage{microtype}
\usepackage{amsmath,amssymb,mathtools}
\usepackage{geometry}
\usepackage{xcolor}
\usepackage{enumitem}
\usepackage{fancyhdr}
\usepackage[hidelinks]{hyperref}

\geometry{left=2.3cm,right=2.3cm,top=2.2cm,bottom=2.25cm,headheight=14pt}
\setlength{\parindent}{0pt}
\setlength{\parskip}{0.42em}
\setlist[enumerate]{topsep=0.3em,itemsep=0.25em,leftmargin=1.8em}
\allowdisplaybreaks

\definecolor{azul}{RGB}{34,73,110}
\definecolor{gris}{RGB}{90,96,102}

\pagestyle{fancy}
\fancyhf{}
\fancyhead[L]{\small\color{gris}GC2D: midpoint ABBA method}
\fancyhead[R]{\small\color{gris}Jacobian calculation}
\fancyfoot[C]{\thepage}
\renewcommand{\headrulewidth}{0.3pt}

\newcommand{\R}{\mathbb{R}}
\newcommand{\Id}{I}
\newcommand{\Y}{\mathcal{Y}}
\newcommand{\M}{\mathcal{M}}
\newcommand{\PhiAB}{\Phi^{\mathrm{ABBA}}_{h,t_n}}

\begin{document}

\begin{center}
  {\LARGE\bfseries\color{azul}Jacobian of midpoint ABBA with mean projection}\par
  \vspace{0.35em}
  {\large Calculation for the GC2D problem}\par
  \vspace{0.4em}
  {\small Midpoint method without multipliers or Newton iterations}
\end{center}

\vspace{0.45em}

\section{Sistema físico y notación}

El estado físico es
\[
 y=(X,Y)\in\R^2,
\]
y satisface
\[
 \dot y=f(y,t),
 \qquad
 f(y,t)=J_0\nabla_y\phi(y,t),
 \qquad
 J_0=
 \begin{pmatrix}
  0&-1\\
  1&0
 \end{pmatrix}.
\]
Por tanto,
\[
 \dot X=-\partial_Y\phi(t,X,Y),
 \qquad
 \dot Y=\partial_X\phi(t,X,Y).
\]

Denotaremos por
\[
 W(y,t):=D_yf(y,t)=J_0\nabla_y^2\phi(y,t)
\]
la jacobiana del campo vectorial. En coordenadas,
\[
 \boxed{
 W(y,t)=
 \begin{pmatrix}
  -\phi_{XY}(y,t)&-\phi_{YY}(y,t)\\
  \phi_{XX}(y,t)&\phi_{XY}(y,t)
 \end{pmatrix}.}
 \tag{1}
\]

En todo el cálculo, \(t_n\) y \(h\) se consideran fijos. La derivación se realiza
únicamente con respecto al estado inicial \(y_n\).

\section{Definición del paso ABBA explícito}

Introducimos dos copias del estado físico y comenzamos cada paso sobre la diagonal:
\[
 u^{(0)}=v^{(0)}=y_n.
 \tag{2}
\]
Tomando
\[
 s=\frac h2,
 \qquad
 t_{n+1}=t_n+h,
\]
las cuatro etapas ABBA son
\begin{align}
 u^{(1)}
 &=u^{(0)}+s f\bigl(v^{(0)},t_n\bigr),
 \tag{3a}\\
 v^{(1)}
 &=v^{(0)}+s f\bigl(u^{(1)},t_n\bigr),
 \tag{3b}\\
 v^{(2)}
 &=v^{(1)}+s f\bigl(u^{(1)},t_{n+1}\bigr),
 \tag{3c}\\
 u^{(2)}
 &=u^{(1)}+s f\bigl(v^{(2)},t_{n+1}\bigr).
 \tag{3d}
\end{align}
El nuevo estado físico es la media de las dos copias finales:
\[
 \boxed{
 y_{n+1}=\Psi_{h,t_n}(y_n)
 =\frac{u^{(2)}+v^{(2)}}2.}
 \tag{4}
\]

Las matrices \(W\) que aparecen al derivar deben evaluarse en las etapas concretas
del método. Definimos
\begin{align*}
 W_1&:=W\bigl(v^{(0)},t_n\bigr),
 &W_2&:=W\bigl(u^{(1)},t_n\bigr),\\
 W_3&:=W\bigl(u^{(1)},t_{n+1}\bigr),
 &W_4&:=W\bigl(v^{(2)},t_{n+1}\bigr),
\end{align*}
y agrupamos las dos etapas centrales mediante
\[
 \boxed{S:=W_2+W_3.}
 \tag{5}
\]
Todas estas matrices son de tamaño \(2\times2\).

\section{Propagación de las derivadas por las etapas}

Para hacer visible la regla de la cadena, introducimos
\[
 U_0:=D_{y_n}u^{(0)},\quad
 V_0:=D_{y_n}v^{(0)},\quad
 U_1:=D_{y_n}u^{(1)},
\]
\[
 V_1:=D_{y_n}v^{(1)},\quad
 V_2:=D_{y_n}v^{(2)},\quad
 U_2:=D_{y_n}u^{(2)}.
\]
Como las dos copias iniciales son exactamente \(y_n\), se tiene
\[
 \boxed{U_0=V_0=\Id_2.}
 \tag{6}
\]

\subsection{Primera etapa}

Derivando (3a),
\[
 U_1
 =U_0+sW_1V_0.
\]
Usando (6), obtenemos
\[
 \boxed{U_1=\Id_2+sW_1.}
 \tag{7}
\]

\subsection{Segunda etapa}

De (3b), por la regla de la cadena,
\[
 V_1=V_0+sW_2U_1.
\]
Sustituyendo (6) y (7),
\begin{align*}
 V_1
 &=\Id_2+sW_2(\Id_2+sW_1)\\
 &=\Id_2+sW_2+s^2W_2W_1.
\end{align*}
Por tanto,
\[
 \boxed{V_1=\Id_2+sW_2+s^2W_2W_1.}
 \tag{8}
\]

\subsection{Tercera etapa}

La etapa (3c) utiliza de nuevo \(u^{(1)}\), pero evalúa el campo en
\(t_{n+1}\). Por tanto,
\[
 V_2=V_1+sW_3U_1.
\]
Sustituyendo (7) y (8),
\begin{align*}
 V_2
 &=\Id_2+sW_2+s^2W_2W_1
   +sW_3(\Id_2+sW_1)\\
 &=\Id_2+s(W_2+W_3)+s^2(W_2+W_3)W_1.
\end{align*}
Con la notación \(S=W_2+W_3\), resulta
\[
 \boxed{V_2=\Id_2+sS+s^2SW_1.}
 \tag{9}
\]

\subsection{Cuarta etapa}

Finalmente, al derivar (3d),
\[
 U_2=U_1+sW_4V_2.
\]
Sustituyendo (7) y (9),
\begin{align*}
 U_2
 &=\Id_2+sW_1
   +sW_4\bigl(\Id_2+sS+s^2SW_1\bigr)\\
 &=\Id_2+s(W_1+W_4)
   +s^2W_4S+s^3W_4SW_1.
\end{align*}
Por tanto,
\[
 \boxed{
 U_2
 =\Id_2+s(W_1+W_4)
 +s^2W_4S+s^3W_4SW_1.}
 \tag{10}
\]

Es importante respetar el orden de los productos: las matrices \(W_i\), evaluadas
en puntos distintos, no conmutan en general.

\section{Jacobiana del mapa físico}

Derivando la media final (4), se obtiene directamente
\[
 D\Psi_{h,t_n}(y_n)
 =\frac12\left(U_2+V_2\right).
 \tag{11}
\]
Sustituyendo (9) y (10),
\begin{align*}
 D\Psi_{h,t_n}(y_n)
 &=\frac12\Bigl[
  \Id_2+s(W_1+W_4)+s^2W_4S+s^3W_4SW_1\\
 &\hspace{3.5cm}
  +\Id_2+sS+s^2SW_1
 \Bigr].
\end{align*}
Agrupando términos del mismo orden en \(s\), llegamos a la fórmula exacta
\[
 \boxed{
 \begin{aligned}
 D\Psi_{h,t_n}(y_n)
 ={}&\Id_2
 +\frac{s}{2}(W_1+S+W_4)\\
 &+\frac{s^2}{2}\bigl(W_4S+SW_1\bigr)
 +\frac{s^3}{2}W_4SW_1.
 \end{aligned}}
 \tag{12}
\]
Como \(S=W_2+W_3\), la misma expresión completamente desarrollada es
\[
 \boxed{
 \begin{aligned}
 D\Psi_{h,t_n}(y_n)
 ={}&\Id_2
 +\frac{s}{2}(W_1+W_2+W_3+W_4)\\
 &+\frac{s^2}{2}
 \left[W_4(W_2+W_3)+(W_2+W_3)W_1\right]\\
 &+\frac{s^3}{2}W_4(W_2+W_3)W_1.
 \end{aligned}}
 \tag{13}
\]
Finalmente, usando \(s=h/2\),
\[
 \boxed{
 \begin{aligned}
 D\Psi_{h,t_n}(y_n)
 ={}&\Id_2
 +\frac{h}{4}(W_1+W_2+W_3+W_4)\\
 &+\frac{h^2}{8}
 \left[W_4(W_2+W_3)+(W_2+W_3)W_1\right]\\
 &+\frac{h^3}{16}W_4(W_2+W_3)W_1.
 \end{aligned}}
 \tag{14}
\]
Esta es la jacobiana \(2\times2\) del paso físico
\(y_n\mapsto y_{n+1}\).

\section{Deducción equivalente mediante la jacobiana extendida}

También puede obtenerse (12) a partir del mapa ABBA en el espacio duplicado.
Escribimos
\[
 D\PhiAB=
 \begin{pmatrix}
  A&B\\
  C&D
 \end{pmatrix},
 \tag{15}
\]
donde cada bloque es \(2\times2\). La derivación de las etapas con respecto a
\(u^{(0)}\) y \(v^{(0)}\) proporciona
\[
 \boxed{
 \begin{aligned}
 A&=\Id_2+s^2W_4S,\\
 B&=s(W_1+W_4)+s^3W_4SW_1,\\
 C&=sS,\\
 D&=\Id_2+s^2SW_1.
 \end{aligned}}
 \tag{16}
\]

La inclusión diagonal y la extracción de la media son
\[
 E=
 \begin{pmatrix}
  \Id_2\\ \Id_2
 \end{pmatrix},
 \qquad
 P=\frac12
 \begin{pmatrix}
  \Id_2&\Id_2
 \end{pmatrix}.
 \tag{17}
\]
Como
\[
 \Psi_{h,t_n}=P\PhiAB E,
\]
la regla de la cadena da
\begin{align*}
 D\Psi_{h,t_n}(y_n)
 &=P(D\PhiAB)E\\
 &=\frac12
 \begin{pmatrix}
  \Id_2&\Id_2
 \end{pmatrix}
 \begin{pmatrix}
  A&B\\
  C&D
 \end{pmatrix}
 \begin{pmatrix}
  \Id_2\\ \Id_2
 \end{pmatrix}\\
 &=\boxed{\frac12(A+B+C+D)}.
 \tag{18}
\end{align*}
Al sustituir los cuatro bloques de (16), se recupera exactamente (12).

\section{Forma entrada por entrada}

Definimos la matriz de corrección
\[
 \mathcal T
 :=\frac{s}{2}(W_1+S+W_4)
 +\frac{s^2}{2}(W_4S+SW_1)
 +\frac{s^3}{2}W_4SW_1
 =:(\tau_{ij})_{1\leq i,j\leq2}.
 \tag{19}
\]
Entonces
\[
 \boxed{
 D\Psi_{h,t_n}(y_n)
 =
 \begin{pmatrix}
  1+\tau_{11}&\tau_{12}\\
  \tau_{21}&1+\tau_{22}
 \end{pmatrix}.}
 \tag{20}
\]
Por definición, esta matriz es
\[
 D\Psi_{h,t_n}(y_n)
 =
 \begin{pmatrix}
  \dfrac{\partial X_{n+1}}{\partial X_n}
  &
  \dfrac{\partial X_{n+1}}{\partial Y_n}\\[2.2ex]
  \dfrac{\partial Y_{n+1}}{\partial X_n}
  &
  \dfrac{\partial Y_{n+1}}{\partial Y_n}
 \end{pmatrix}.
 \tag{21}
\]
Las entradas \(\tau_{ij}\) se calculan mediante productos ordinarios de matrices
\(2\times2\), usando para cada \(W_i\) la expresión de las segundas derivadas de
\(\phi\) dada en (1).

\section{Comprobación con un potencial cuadrático}

Consideremos
\[
 \phi(X,Y)=\frac12(X^2+Y^2).
\]
Entonces
\[
 f(y)=J_0y,
 \qquad
 W_1=W_2=W_3=W_4=J_0,
 \qquad
 S=2J_0,
\]
y \(J_0^2=-\Id_2\). Sustituyendo en (12),
\begin{align*}
 D\Psi_h
 &=\Id_2+2sJ_0-2s^2\Id_2-s^3J_0\\
 &=\left(1-\frac{h^2}{2}\right)\Id_2
   +\left(h-\frac{h^3}{8}\right)J_0.
\end{align*}
Por tanto,
\[
 \boxed{
 D\Psi_h=
 \begin{pmatrix}
  1-\dfrac{h^2}{2}&-h+\dfrac{h^3}{8}\\[1.2ex]
  h-\dfrac{h^3}{8}&1-\dfrac{h^2}{2}
 \end{pmatrix}.}
 \tag{22}
\]
Su determinante es
\[
 \det(D\Psi_h)
 =\left(1-\frac{h^2}{2}\right)^2
  +\left(h-\frac{h^3}{8}\right)^2
 =1+\frac{h^6}{64}.
 \tag{23}
\]
Así, para \(h\neq0\), la jacobiana no es simpléctica. Esta comprobación coincide
con el mapa lineal obtenido directamente para el mismo ejemplo y confirma tanto la
fórmula (12) como la limitación geométrica del promedio final.

\section{Cálculo práctico en cada paso}

Para evaluar la jacobiana junto con un paso del método:
\begin{enumerate}
 \item Calcular las cuatro etapas (3a)--(3d).
 \item Evaluar \(W_1,W_2,W_3,W_4\) en los puntos y tiempos indicados.
 \item Formar \(S=W_2+W_3\).
 \item Calcular los productos \(W_4S\), \(SW_1\) y \(W_4SW_1\).
 \item Construir \(D\Psi_{h,t_n}(y_n)\) mediante (12) o (14).
\end{enumerate}

No es necesario resolver ningún sistema lineal ni invertir ninguna matriz. Tampoco
aparecen derivadas terceras de \(\phi\): basta con evaluar la hessiana de \(\phi\) en
las cuatro etapas para construir las matrices \(W_i\).

\section{Resultado final}

La jacobiana exacta del método ABBA explícito con proyección mediante la media es
\[
 \boxed{
 D\Psi_{h,t_n}(y_n)
 =\Id_2
 +\frac{s}{2}(W_1+S+W_4)
 +\frac{s^2}{2}(W_4S+SW_1)
 +\frac{s^3}{2}W_4SW_1,
 \qquad
 S=W_2+W_3.}
\]
Esta matriz describe la sensibilidad del nuevo estado \(y_{n+1}\) respecto del estado
inicial \(y_n\). A diferencia del método con proyección implícita, aquí no aparece la
derivada de ningún multiplicador: toda la dependencia se propaga explícitamente por las
cuatro etapas y por la media final.

\end{document}
